\begin{thebibliography}{99}

\bibitem[Banco Mundial(2024)]{wb2024} Banco Mundial (2024). Memorando Económico del País de Colombia: Crecer para prosperar. \emph{Banco Mundial, Washington DC}. \href{https://documents1.worldbank.org/curated/en/099042525165519971/pdf/P177730-ca16ae09-fb9a-4f6a-a7f4-b88f1347d370.pdf}{Fuente}.

\bibitem[Bartelsman et al.(2013)]{bartelsman2013} Bartelsman, Eric; Haltiwanger, John; Scarpetta, Stefano (2013). Cross-Country Differences in Productivity: The Role of Allocation and Selection. \emph{American Economic Review, 103(1), 305-334}. \href{https://www.aeaweb.org/articles?id=10.1257/aer.103.1.305}{Fuente}.

\bibitem[Ben Hassine(2019)]{benhassine2019} Ben Hassine, Haithem (2019). Productivity Growth and Resource Reallocation in France: The Process of Creative Destruction. \emph{Economie et Statistique / Economics and Statistics, 507-508, 115-133}. \href{https://doi.org/10.24187/ecostat.2019.507d.1979}{Fuente}.

\bibitem[Brown et al.(2018)]{brown2018} Brown, J. David; Crespi, Gustavo A.; Iacovone, Leonardo; Marcolin, Luca (2018). Decomposing firm-level productivity growth and assessing its determinants: evidence from the Americas. \emph{The Journal of Technology Transfer, 43, 1571-1606}. \href{https://doi.org/10.1007/s10961-018-9678-0}{Fuente}.

\bibitem[César y Falcone(2020)]{cesar2020} César, Andrés; Falcone, Guillermo (2020). Heterogeneous Effects of Chinese Import Competition on Chilean Manufacturing Plants. \emph{Economía, 20(2), 1-60}. \href{https://www.jstor.org/stable/27007012}{Fuente}.

\bibitem[Eslava et al.(2004)]{eslava2004} Eslava, Marcela; Haltiwanger, John; Kugler, Adriana; Kugler, Maurice (2004). The effects of structural reforms on productivity and profitability enhancing reallocation: evidence from Colombia. \emph{Journal of Development Economics, 75(2), 333-371}. \href{https://www.nber.org/papers/w10367}{Fuente}.

\bibitem[Eslava et al.(2006)]{eslava2006} Eslava, Marcela; Haltiwanger, John; Kugler, Adriana; Kugler, Maurice (2006). Plant Turnover and Structural Reforms in Colombia. \emph{IMF Staff Papers, 53, número especial}. \href{https://www.imf.org/external/pubs/ft/staffp/2006/03/eslava.htm}{Fuente}.

\bibitem[Foster et al.(2001)]{foster2001} Foster, Lucia; Haltiwanger, John; Krizan, C. J. (2001). Aggregate Productivity Growth: Lessons from Microeconomic Evidence. \emph{New Developments in Productivity Analysis, University of Chicago Press, capítulo 8}. \href{https://www.nber.org/books-and-chapters/new-developments-productivity-analysis/aggregate-productivity-growth-lessons-microeconomic-evidence}{Fuente}.

\bibitem[Griliches y Regev(1995)]{griliches1995} Griliches, Zvi; Regev, Haim (1995). Firm productivity in Israeli industry 1979-1988. \emph{Journal of Econometrics, 65(1), 175-203}. \href{https://www.sciencedirect.com/science/article/pii/030440769401601U}{Fuente}.

\bibitem[Harris y Moffat(2019)]{harris2019} Harris, Richard; Moffat, John (2019). The Decline of British Manufacturing, 1973-2012: The Role of Total Factor Productivity. \emph{National Institute Economic Review, 247(1), R19-R31}. \href{https://doi.org/10.1177/002795011924700112}{Fuente}.

\bibitem[Iacovone et al.(2021)]{mexico2021} Iacovone, Leonardo; Muñoz Moreno, Rafael; Olaberria, Eduardo; Pereira López, Mariana de la Paz (2021). Productivity Growth in Mexico: Understanding Main Dynamics and Key Drivers. \emph{Banco Mundial; edición PDF con copyright 2022, cita recomendada 2021}. \href{https://documents1.worldbank.org/curated/en/099100103252224294/pdf/P17082908942250380867901fba9060dcbc.pdf}{Fuente}.

\bibitem[IFC(2022)]{ifc2022} International Finance Corporation (2022). Creating Markets in Chile: A Stronger Private Sector for a More Productive and Inclusive Society. \emph{Country Private Sector Diagnostic, junio 2022}. \href{https://www.ifc.org/content/dam/ifc/doc/mgrt/cpsd-chile.pdf}{Fuente}.

\bibitem[FMI(2025)]{imf2025} International Monetary Fund (2025). Fostering Growth through Business Dynamism. \emph{Regional Economic Outlook: Western Hemisphere, octubre 2025, capítulo 2}. \href{https://www.imf.org/-/media/files/publications/reo/whd/2025/october/english/ch2.pdf}{Fuente}.

\bibitem[FMI(2024)]{colombia2024} International Monetary Fund (2024). Colombia: Selected Issues. \emph{IMF Country Report 24/83; capítulo sobre productividad}. \href{https://www.elibrary.imf.org/view/journals/002/2024/083/article-A001-en.xml}{Fuente}.

\bibitem[FMI(2022)]{portugal2022} International Monetary Fund (2022). Portugal: Selected Issues. \emph{IMF Country Report 22/204; capítulo sobre productividad}. \href{https://www.elibrary.imf.org/view/journals/002/2022/204/article-A001-en.xml}{Fuente}.

\bibitem[Melitz y Polanec(2015)]{melitz2015} Melitz, Marc J.; Polanec, Sašo (2015). Dynamic Olley-Pakes productivity decomposition with entry and exit. \emph{RAND Journal of Economics, 46(2), 362-375}. \href{https://doi.org/10.1111/1756-2171.12088}{Fuente}.

\bibitem[Olley y Pakes(1996)]{olley1996} Olley, G. Steven; Pakes, Ariel (1996). The Dynamics of Productivity in the Telecommunications Equipment Industry. \emph{Econometrica, 64(6), 1263-1297}. \href{https://www.nber.org/papers/w3977}{Fuente}.

\bibitem[Pancost y Yeh(2022)]{pancost2022} Pancost, N. Aaron; Yeh, Chen (2022). Decomposing Aggregate Productivity. \emph{U.S. Census Bureau, CES Working Paper 22-25}. \href{https://www.census.gov/library/working-papers/2022/adrm/CES-WP-22-25.html}{Fuente}.

\bibitem[Pérez et al.(2024)]{opce2024} Pérez García, Francisco; Mas Ivars, Matilde; Pilat, Dirk (dirs.) (2024). La productividad en España: los motores de cambio. Informe OPCE 2024. \emph{Fundación BBVA-Ivie, versión preliminar de noviembre de 2024}. \href{https://www.fbbva.es/wp-content/uploads/2025/01/DE_2025_productividad_espana_motores_de_cambio_informe_opce_2024_version_preliminar.pdf}{Fuente}.

\bibitem[Song(2021)]{song2021} Song, Yeongkwan (2021). Trade Liberalization and Manufacturing Productivity Changes in Korea during the Past Three Decades. \emph{KDI Journal of Economic Policy, 43(1), 53-80}. \href{https://doi.org/10.23895/kdijep.2021.43.1.53}{Fuente}.

\end{thebibliography}